\documentclass[11pt,a4paper]{article}
\usepackage[utf8]{inputenc}
\usepackage[T1]{fontenc}
\usepackage{hyperref}
\usepackage{xcolor}
\usepackage{listings}
\usepackage{graphicx}
\usepackage{geometry}
\usepackage{titlesec}
\usepackage{enumitem}
\usepackage{fancyvrb}
\usepackage{bookmark}

% ── DESIGN TOKENS ─────────────────────────────────────────────────────────────
\definecolor{primaryblue}{HTML}{0071e3}
\definecolor{darkbg}{HTML}{050508}
\definecolor{accentpurple}{HTML}{bf5af2}
\definecolor{successgreen}{HTML}{30d158}

\geometry{margin=0.85in}
\hypersetup{
    colorlinks=true,
    linkcolor=primaryblue,
    filecolor=magenta,      
    urlcolor=primaryblue,
}

\lstset{
    basicstyle=\ttfamily\footnotesize,
    breaklines=true,
    frame=single,
    backgroundcolor=\color{gray!2},
    keywordstyle=\color{primaryblue}\bfseries,
    commentstyle=\color{gray},
    stringstyle=\color{successgreen},
    showstringspaces=false
}

\titleformat{\section}{\Large\bfseries\color{primaryblue}}{\thesection}{1em}{}[{\titlerule[0.5pt]}]
\titleformat{\subsection}{\large\bfseries}{\thesubsection}{1em}{}

\title{\textbf{SEO Insight: Master Technical Documentation}}
\author{\textbf{Aditya} (Lead Developer) \& \textbf{Antigravity AI}}
\date{\today}

\begin{document}

\maketitle
\begin{abstract}
    This document serves as the comprehensive technical manual for SEO Insight, a high-performance, local-first SEO intelligence platform. It covers the architectural design, AI integration strategies, professional UI systems, and the roadmap for commercial scalability.
\end{abstract}

\newpage
\tableofcontents
\newpage

% ── Section 1: Overview ──────────────────────────────────────────────
\section{Project Definition \& Philosophy}
SEO Insight is built to challenge the current SEO SaaS landscape by providing enterprise-grade scraping and AI analysis in a private, local-first environment. 

\subsection{Architectural Pillars}
\begin{enumerate}
    \item \textbf{Local-First Persistence}: Uses SQLite for zero-latency data access.
    \item \textbf{AI Velocity}: Leverages Groq (Llama 3.3) for sub-second inference.
    \item \textbf{Visual Excellence}: A premium Apple-inspired design system.
\end{enumerate}

% ── Section 2: Setup ──────────────────────────────────────────────────
\section{Setup \& Orchestration}
The project uses a hybrid setup model to ensure ease of use for developers and robustness for production environments.

\subsection{Automated Setup (setup.py)}
The Python-based setup wizard handles:
\begin{itemize}
    \item Dependency verification (Node.js, pnpm).
    \item Environment variable validation.
    \item Automated package installation and dev-server launch.
\end{itemize}

% ── Section 3: Tech Stack Deep Dive ──────────────────────────────────
\section{Technical Stack Deep Dive}

\subsection{Frontend: Next.js 15}
Utilizes Turbopack for near-instant development cycles and Server Components to minimize client-side JavaScript overhead.

\subsection{AI: Groq & Perplexity}
\begin{itemize}
    \item \textbf{Groq}: Primary engine for report synthesis and global chat.
    \item \textbf{Perplexity}: Optional investigator module for deep entity verification and backlink discovery.
\end{itemize}

\subsection{Scraping: Bright Data}
Manages the SERP Snapshot lifecycle, bypassing advanced bot detections to provide raw search data.

% ── Section 4: The Analysis Pipeline ─────────────────────────────────
\section{The Analysis Pipeline (State Machine)}
The generation of an SEO report follows a strict 6-phase state machine:
\begin{enumerate}
    \item \textbf{Trigger}: Prompt ingestion and initial DB job creation.
    \item \textbf{Scraping}: Snapshot creation and async polling.
    \item \textbf{Compression}: Trimming raw HTML to essential tokens.
    \item \textbf{Inference}: AI reasoning using Llama 3.3.
    \item \textbf{Hardening}: Zod-based data validation and repair.
    \item \textbf{Persistence}: Final JSON storage and UI state update.
\end{enumerate}

% ── Section 5: Intelligence Modules ──────────────────────────────────
\section{Intelligence Modules}

\subsection{Aditya Mode (Developer Mode)}
Triggered by the phrase "This is Aditya", this mode unlocks the internal architecture knowledge base for the AI, allowing it to assist in project maintenance.

\subsection{Deep Research (Perplexity)}
An investigative layer that cross-references LinkedIn, GitHub, and professional registries to ensure entity accuracy.

% ── Section 6: Design System ──────────────────────────────────────────
\section{The Design System}
The UI is built on a series of professional CSS and React tokens.

\subsection{Visual Tokens}
\begin{itemize}
    \item \textbf{Aurora Background}: Floating blurred orbs with motion-driven drift.
    \item \textbf{Glassmorphism}: 24px blur radius with semi-transparent borders.
    \item \textbf{Staggered Entrance}: 20ms incremental delays for smoother content delivery.
\end{itemize}

\subsection{Page Progress}
A custom gradient loader that simulates organic growth during route transitions, ensuring the UI always feels responsive.

% ── Section 7: Scalability & Security ────────────────────────────────
\section{Scalability \& Security}

\subsection{Middleware (The Gatekeeper)}
Currently set to passthrough mode for local development. It is primed to integrate with \textbf{Clerk} for route-level protection in a multi-user environment.

\subsection{Commercial Ready (Pricing)}
The platform includes a pre-designed pricing structure (\texttt{app/pricing/page.tsx}) ready for Stripe integration.

% ── Section 8: Troubleshooting ──────────────────────────────────────
\section{Maintenance \& Troubleshooting}
\begin{itemize}
    \item \textbf{Rate Limits}: Managed via the \texttt{AnalysisLock} singleton.
    \item \textbf{Scraping Failures}: Handled via the UI-integrated "Retry" engine.
    \item \textbf{Cancellation}: Active jobs can be halted using the "Stop Generation" API.
\end{itemize}

\end{document}
